\documentclass[11pt]{article}
\usepackage[czech]{babel}
\usepackage[a4paper, left=2.0cm,top=3.0cm,textheight=24.0cm,textwidth=17.0cm]{geometry}
\usepackage[utf8]{inputenc}
\usepackage{times}
\usepackage{csquotes}
\usepackage{marvosym}
\usepackage{url}
\usepackage[hidelinks,unicode]{hyperref}
\usepackage{graphicx}
\def\UrlBreaks{\do\/\do-}
\Urlmuskip=0mu plus 1mu

\graphicspath{ {./pics/} }

\hypersetup{
    linktocpage,
    allcolors=blue,
    colorlinks=true, %set true if you want colored links
    linktoc=all,     %set to all if you want both sections and subsections linked
    linkcolor=blue,  %choose some color if you want links to stand out
    citecolor=blue,
}

\setlength{\parindent}{1em}

\begin{document}
\selectlanguage{czech}
\begin{titlepage}
    \begin{center}    
        \textsc{\Huge{Vysoké učení technické v~Brně} \\[0.4em] \huge{Fakulta informačních technologií} }\\
        \vspace{\stretch{0.382}}
        \LARGE{Síťové aplikace a správa sítí}\\
        \Huge{Implementace TFTP klienta}
        \vspace{\stretch{0.618}}  
    \end{center}
    {\Large\hfill Ondřej Keprt}\\
    {\Large \today \hfill \href{mailto:xkeprt03@stud.fit.vutbr.cz}{xkeprt03@stud.fit.vutbr.cz}}
\end{titlepage}

\newpage
\tableofcontents
\thispagestyle{empty}
\newpage
\setcounter{page}{1}

\section{Úvod}
Cílem projektu bylo implementovat Trivial File Transfer Protocol (TFTP) klienta. TFTP je jednoduchý protokol pro stahovaní dat serveru nebo ukládání dat na server.
Implementovaný klient nepodporuje mód přenosu \texttt{mail}.

\section{Protokol}
\subsection{Zakladní verze}
Základní verze je specifikováná v \texttt{RFC~1350}. \cite{revision2}

Protokol pracuje s UDP datagramy, tedy s nespolehlivým doručováním paketů. Proto je protokol opatřen vlastním mechanismem pro potvrzení dat.
Klient zasílá požadavek na server ve kterém specifikuje, zda chce data číst nebo nahrávat. Na první místo v paketu tedy uloží dvou bytovou 
hodnotu 1 (čtení) nebo 2 (zápis). Další položkou je řetězec znaků, ukončených bytem s hodnotou 0, který specifikuje název souboru. Další položkou paketu
je pak řetězec znaků , který určuje mód přenosu (binary/octet nebo ascii/netascii) \ref{pict:RFC_1350_packets}.

Server na to odpovídá datovým paketem, kde na prvním místě je dvoubytová hodnota specifikující data paket~(3), následuje dvou bytová hodnota,
která specifikuje blok dat. Poté následuje n\footnote{Základní verze udává velikost 512 bytů} bytů. Pokud je počet dat menší jak n, jedná se o poslední data transportu \ref{pict:RFC_1350_packets}.

Uživatel i server na data odpovidají potvrzovacím paketem, který má dvou bytovou hodnotu pro potvrzení~(4). Následuje dvou bytová hodnota, která určuje přijatý blok dat \ref{pict:RFC_1350_packets}.

Pokud nastane nějaká chyba, je zaslán paket, který značí chybu. Obsahuje dvou bytovou hodnotu, specifikující chybový paket~(5), dvou bytovou hodnotu pro specifikaci
chyby a následuje řetězec znaků, obsahující chybovou hlášku \ref{pict:RFC_1350_packets}.  

\begin{figure}[ht]
    \label{pict:RFC_1350_packets} 
    \centering
    \includegraphics[scale=2.0]{RFC_1350_packets.eps}
    \caption{Data v paketech ve verzi 2 bez rozšiřujících voleb \cite{revision2}}        
\end{figure}

\subsection{Rozšířená verze}
Na základní verzi navazuji rozšíření, která upravují parametry přenosu dat.

V \texttt{RFC 2347} je specifikován způsob zasílání a potvrzování voleb. Za packet ze základní verze přidávám postupně volby a to způsobem řetězec znaků,
který specifikuje volbu. Následuje řetězec znaků, který specifikuje hodnotu volby. Pokud chceme více voleb najednou, přidáváme nadále
na konec paketu vždy specifikaci volby a její hodnotu, stejně jako při první volbě \ref{pict:RFC_1347_option_list.eps}. 

\begin{figure}[ht]
    \label{pict:RFC_1347_option_list.eps} 
    \centering
    \includegraphics[scale=2.0]{RFC_2347_option_list.eps}
    \caption{Data v paketech v rozšíření \cite{option_extension}}        
\end{figure}

Na tento paket server odpovídá paketem, který obsahuje dvou bytovou hodnotu pro potvrzení přijetí voleb (6), následovaný žetězcem znaků pro specifikaci znaků a řetězcem pro její hodnotu.
Pokud uživatel, chce číst ze serveru, tento paket musí potvrdit datový paketem, jako hodnotu bloku uvede 0. Server poté začne posílat/přijímat datové pakety. Pokud uživatel chce zapisovat na server, odesílá rovnou první datový paket \ref{pict:RFC_1347_option_list.eps}. \cite{option_extension}

\begin{figure}[ht]
    \label{pict:RFC_1347_option_list_ack.eps} 
    \centering
    \includegraphics[scale=2.0]{RFC_2347_option_list_ack.eps}
    \caption{Paket potvrzující rozšíření \cite{option_extension}}        
\end{figure}

Jedna z voleb je \texttt{tsize}, které specifikuje velikost přenášeného souboru \ref{pict:RFC_1349_tsize.eps}.

\begin{figure}[ht]
    \label{pict:RFC_1349_tsize.eps} 
    \centering
    \includegraphics[scale=2.0]{RFC_2349_tsize.eps}
    \caption{Paket pro volbu tsize \cite{extension_timeout_tsize}}        
\end{figure}

Další volbou je \texttt{timeout}, která specifikuje dobu, po kterou server čeká před znovu poslání nepotvrzených dat.
Pokud server odmítne toto nastavení, v odpovědi tuto volbu neuvede \ref{pict:RFC_1349_timeout.eps}.

\begin{figure}[ht]
    \label{pict:RFC_1349_timeout.eps} 
    \centering
    \includegraphics[scale=2.0]{RFC_2349_timeoUt.eps}
    \caption{Paket pro volbu timeout \cite{extension_timeout_tsize}}        
\end{figure}

Rozšíření specifikující velikost bloku dat přenášených dat nese název \texttt{blksize}. Pokud server odmítne tuto hodnotu,
v odpovědi specifikuje největší akceptovatelnou velikost, která je menší, než kterou uživatel navrhoval \ref{pict:RFC_1348_blksize.eps}.

\begin{figure}[ht]
    \label{pict:RFC_1348_blksize.eps} 
    \centering
    \includegraphics[scale=2.0]{RFC_2348_blksize}
    \caption{Paket pro volbu timeout \cite{extension_blksize}}        
\end{figure}

\subsection{Výhody a nevýhody}
Jde o jednoduchý protokol, který není náročný na implementaci a tedy na místo v paměti. Typicky se využíval pro bootování 
operačního systému na bezdiskových strojích ze serveru v místní síti. \cite{WIKI}

Protokol se nehodí pro komunikaci mimo lokální síť protože není nijak zabezpečený.

Protokol se nehodí pro sítě, které mají velkou latenci. Jelikož protokol posílá jeden paket a čeká, až protistrana potvrdí jeho přijetí. 
Najednou se tedy posílá pouze jeden paket a to prodlužuje dobu přenosu. \cite{WIKI}
 
\section{Implementace}
Program byl testován pomocí serveru \href{https://pjo2.github.io/tftpd64/}{TFTPD32}. Podrobný popis implementace a postup programu je popsán přímo ve zdrojových kódech programu.

\subsection{Příkazová řádka}
Klient, kterého jsem vytvořil obsahuje příkazovou řádku, na kterou se zadávají parametry pro přenos. Tato příkazová řádka je blokující, tedy
další příkaz může být zadán až po ukončení přenosu. Pomocí příkazy \texttt{-h} se zobrazí nápověda, jaké parametry lze možné zadávat. Pokud uživatel
zadá nesprávné parametry, program vypíše chybovou hlášku a bude očekávat další vstup. Program lze ukončit napsaním \texttt{end}.

\subsection{Reakce na volby}
Pokud server nepodporoval volby, uživatel byl o tomto obeznámen a přenos dat pokračoval dál podle znění \texttt{RFC~1350}. ~\cite{revision2}
Tady datový blok měl velikost 512 bytů a žádné další volby nebyly aplikovány.

Při každém odesílání požadavku se posílala volba \texttt{tsize}, což je velikost souboru. Pokud server při čtení odmítl sdělit velikost souboru a začal 
odesílat rovnou tato data byla tato data přijata. Pokud server vrátil velikost požadovaného souboru, ověřilo se volné místo na disku, při dostatku
volného místa se ze pokračovalo v přenosu, opačně se odeslal chybový paket s příslušným kódem pro nedostatek paměti.

Pokud server při volbě \texttt{blksize} odpověděl menší velikostí datového bloku, přenos pokračoval přenosem těchto menších bloků a uživatel
byl na výstupu obeznámen s touto skutečností.

Pokud uživatel zadal volbu \texttt{timeout} byla tato časová doba okamžitě nastavena pro přenos. Pokud server odmínul tudo dobu, \texttt{timeout}
klienta byl nastavena na základní dobu (5 sekund).

\subsection{Neimplementováno}
Klient nepodporuje možnost posílání dat pomocí multicastové adresy. \\
Klient není schopen přizpůsobit velikost bloku dat rozhraní MTU.

\newpage
%http://www.fit.vutbr.cz/~martinek/latex/czechiso.html, soubor byl, mírně upraven
\bibliographystyle{czechiso} 
\renewcommand{\refname}{Literatura}
\bibliography{zdroje}
\end{document}